\documentclass[UTF8,a4paper,11pt,openany]{ctexbook}
\usepackage{../common/statlearnbook}
\SLSetBookMetadata{第 5 册}{统计学习方法讲义 v2.6.0}{采样方法、主题模型与图排序}
\begin{document}
\setcounter{page}{866}
\setcounter{chapter}{35}
\setcounter{section}{2}
\setcounter{figure}{3}
\pagestyle{headings}
\chaptermark{潜在狄利克雷分配}
\sectionmark{折叠 Gibbs 推断}
\textbf{推导第3步：用留一状态的Beta函数比值形成满条件。}
命题把折叠联合分布中所有与候选主题无关的因子约去，再对$K$个正权重统一归一化；这一“先留一、后作比”的顺序也是实现计数守恒的依据。

图\ref{fig:V5-C06-collapsed-gibbs-counts}把一次更新需要读写的两张计数表并列画出：先从旧主题的文档--主题计数和主题--词计数各减一，再计算所有候选权重，抽取新主题后再各加一。
% v2.6.0 figure UID: FIG-P687-01
% Iteration 03: widen the badge-2-to-document-card whitespace corridor.
\tikzset{slfig-FIG-P687-01/.style={font=\fontsize{9.2pt}{11.0pt}\selectfont,every path/.append style={line cap=round,line join=round}}}
\pgfkeysifdefined{/tikz/alt}{}{\tikzset{alt/.code={}}}
\begin{figure}[htbp]
\centering
\begin{tikzpicture}[slfig-FIG-P687-01,
  >=Stealth,
  every node/.style={font=\fontsize{9.2pt}{11pt}\selectfont,align=center,outer sep=0pt},
  action/.style={draw=SLRuleGray,fill=SLSoftGray,rounded corners=2pt,
    minimum height=12mm,text width=36mm,inner xsep=3mm,inner ysep=2mm},
  evidence/.style={draw=SLMidBlue,fill=SLSoftBlue,rounded corners=2pt,
    minimum height=16mm,text width=44mm,inner xsep=3mm,inner ysep=2mm},
  conditional/.style={draw=SLDeepBlue,fill=white,rounded corners=2pt,
    minimum height=17mm,text width=61mm,inner xsep=3mm,inner ysep=2mm},
  badge/.style={circle,fill=SLDeepBlue,text=white,minimum size=6.4mm,inner sep=0pt,
    font=\fontsize{9.0pt}{10.6pt}\selectfont\bfseries},
  branch/.style={-{Stealth[length=2.1mm,width=1.55mm]},draw=SLMidBlue,line width=.78pt},
  loop/.style={-{Stealth[length=2.2mm,width=1.6mm]},draw=SLDeepBlue,line width=.95pt},
  alt={对象—关系—结论：折叠Gibbs更新构成五步闭环：先移除当前词位，再分别读取带负i上标的文档主题计数和主题词计数，两支证据相乘归一化得到满条件，最后抽样并恢复计数后进入下一词位。}]

\node[action] (remove) at (0,3.00)
  {移除当前词位$i$\\从两张计数表各减一};
\node[badge] at (-2.78,3.50) {1};

\node[evidence] (doc) at (-3.75,1.05)
  {读取文档--主题计数\\
   $\displaystyle n_{mk}^{-i}+\alpha_k$\\[-1pt]
   $\displaystyle n_{m\cdot}^{-i}+\alpha_0$};
\node[badge] at (-6.75,1.48) {2};
\node[evidence] (word) at (3.75,1.05)
  {读取主题--单词计数\\
   $\displaystyle n_{kv}^{-i}+\beta_v$\\[-1pt]
   $\displaystyle n_{k\cdot}^{-i}+\beta_0$};
\node[badge] at (0.68,1.48) {3};

\node[conditional] (full) at (0,-1.00)
  {相乘并对$k$归一化\\[-1pt]
   $\displaystyle p(z_i=k\mid-)
   \propto
   \frac{n_{mk}^{-i}+\alpha_k}{n_{m\cdot}^{-i}+\alpha_0}
   \frac{n_{kv}^{-i}+\beta_v}{n_{k\cdot}^{-i}+\beta_0}$};
\node[badge] at (-4.05,-0.53) {4};

\node[action,draw=SLDeepBlue,fill=SLSoftBlue,text width=41mm] (sample) at (0,-2.85)
  {抽样$k^*$并恢复计数\\把词位$i$加回两张表};
\node[badge] at (-3.03,-2.43) {5};

\draw[branch] (remove.south west) -- (doc.north);
\draw[branch] (remove.south east) -- (word.north);
\draw[branch] (doc.south) -- (full.north west);
\draw[branch] (word.south) -- (full.north east);
\draw[loop] (full.south) -- (sample.north);
\draw[loop] (sample.east) -- (6.65,-2.85) --
  node[midway,right=2.4mm,text width=20mm,align=left,font=\fontsize{8.8pt}{10.5pt}\selectfont,text=SLTextGray]
    {下一词位：\\重新从“移除”开始}
  (6.65,3.00) -- (remove.east);

\node[font=\fontsize{9.0pt}{10.6pt}\selectfont,text=SLTextGray] at (0,-3.90)
  {上标$-i$始终保留到抽样完成，防止当前词位对自身重复计数。};
\end{tikzpicture}
\caption{折叠Gibbs更新先临时移除当前词位，再把文档--主题偏好与主题--单词证据相乘形成满条件分布，最后抽取新主题并恢复计数；上标负$i$防止当前词位对自身产生重复计数}
\label{fig:V5-C06-collapsed-gibbs-counts}
\end{figure}

\FloatBarrier
样本之后可用后验均值平滑估计主题比例与主题--词分布；更稳健的做法是对多个保留样本的预测量或计数平滑估计取平均，而不是只报告最后一次扫描。
\end{document}

